% English translation of checked_pages/esq_119_053.tex; source: cote 119, site 53, PDF page 54, typescript page 51.
% AI-translated by GPT-6 Astra (OpenAI, ChatGPT Pro mode).
\scanpage{119/53}
\mspage{51}

of the themes sketched (assuming there is an ``end''...), and
a hundred years would be little for the richest of them !

Behind the apparent disparity of the themes mentioned here,
an attentive reader will perceive, as I do, a deep unity. This
manifests itself in particular through a common source of inspiration,
the geometry of surfaces, present in all these themes, most
often in the foreground. This source, in relation to my mathematical
``past'', represents a renewal, but by no means a
break. Rather, it points the way toward a new approach
to that still mysterious reality, that of ``\uline{motives}'', which
fascinated me more than any other in the last years of that
past (* ). This fascination has by no means vanished;
rather, it forms part of that of the most burning for me of all
the themes mentioned earlier. But today I am no longer,
as I once was, the willing prisoner of endless tasks,
which so often prevented me from launching into the unknown,
mathematical or otherwise. The time of \uline{tasks} is over for me. If
age has brought me anything, it is to be lighter.

January 1984

(* ) On this subject, see my comments in the ``Thematic Sketch''
of 1972 attached to the present report, under the final heading
``motivic ramblings'',\tnote{TN-053-01}{The French titles are ``Esquisse Thématique'' and ``divagations motiviques''.} (loc. cit. pages 17--18).

\margin{53}
